\subsection{Equal-resource shape trade-offs}
All 449 nominal solver paths meet the inherited acceptance checks. Nominal thickness varies from 0.18919 to 0.20955 mm because arc length is compensated to maintain $\am=0.27$ mm. Figure~\ref{fig:new_front} shows a continuous energy--stress trade-off across the visited shapes. The 16-element sinusoidal reference has $\ua=12.8137$ J\,m$^{-2}$ and $\om=0.51505$. The NURBS energy extreme V01 increases stored potential by 7.79\% with 3.91\% greater stress utilization. The NURBS stress extreme V02 reduces utilization by 10.48\%, at the cost of 22.92\% less potential. These are physically matched shape trade-offs, rather than savings obtained by varying height or material quantity.

The Fourier family provides a useful simpler comparator. Its energy extreme exceeds the sine by only 0.66\%, with 0.09\% greater stress utilization; its stress extreme lowers utilization by 2.49\% while sacrificing 5.81\% of potential. Increased parameter count mainly enlarges the observed NURBS trade-off range. It does not produce a shape that improves both objectives over every simpler reference.

\begin{figure*}[t]
\centering\includegraphics[width=.98\textwidth]{fixed_resource_fronts.pdf}
\caption{Fixed-resource direct-FEM comparison at mesh 16. Left: all accepted nominal evaluations and each family's pooled candidate front; the sinusoid is an explicit simpler reference. Right: median and full seed range of pooled-candidate hypervolume. Each paired algorithm run receives exactly 64 expensive paths and the same initial 16 candidates.}\label{fig:new_front}
\end{figure*}

The paired final hypervolume changes for NSGA-II relative to Sobol are $+1.10\%$, $-4.64\%$ and $-1.06\%$ for seeds 4101--4103. Mean values are 0.53906 and 0.54759, respectively. Thus the tested evolutionary search has no consistent front-quality advantage at this budget. It does require fewer geometric proposals: 1571 in total across its three runs, compared with 6091 for Sobol. Both consume 192 expensive paths. Table~S1 gives each run's counts, indicator and elapsed time; differences in geometric filtering and solver conditioning are included in the timings. The result supports retaining a direct sampling baseline when assessing a more elaborate design pipeline.

\subsection{Which nominal choices survive perturbations?}
The ten shortlisted designs receive all 270 prescribed scenario paths. Every submitted solver path meets the inherited algorithmic acceptance checks, while 39 scenario geometries fail the declared admissibility tests: 32 violate the 0.9 mm radius and seven lose half-period monotonicity. The latter all involve Fourier stress extreme V05, despite its minimum radius remaining above 1.11 mm. The largest reference-envelope overshoot in these rejected geometries is 0.00133 mm. This rejection concerns the declared single-trough shape family, rather than an experimentally measured forming failure.

Five candidates pass all 11 axial scenarios. Joint enrichment removes V09, whose minimum radius decreases from 0.90676 to 0.89750 mm. Four designs remain qualified in all 27 scenarios: the sinusoid V00 and NURBS candidates V01, V07 and V08. The nominal minimum-stress choice V02 and nominal normalized knee V03 are rejected, with scenario minimum radii 0.87860 and 0.88926 mm. A knee chosen only in the nominal objective plane therefore does not automatically provide a defensible design margin.

\begin{figure*}[t]
\centering\includegraphics[width=.98\textwidth]{scenario_qualification.pdf}
\caption{Finite-scenario qualification. Left: nominal mesh-24 responses move to the minimum-potential/maximum-stress coordinates over 27 scenarios; green identifies complete geometric qualification. The two extrema may occur in different scenarios. Right: minimum radius after axial and joint enrichment. V05 also violates the separately checked monotonicity condition.}\label{fig:scenario}
\end{figure*}

All four qualified candidates are nondominated within this finite shortlist. Table~\ref{tab:qualified_main} gives their sampled extrema. V01 retains 7.91\% greater minimum potential than the sine, with 3.38\% greater maximum stress utilization. V08 offers 2.97\% lower maximum utilization, with 12.20\% less minimum potential. V07 trades response for a substantially larger radius margin. The result is a qualified menu of engineering choices, rather than one preference-independent optimum.

\begin{table*}[t]
\centering\small
\caption{Qualified designs over the enriched 27-scenario set, evaluated at mesh 24. Material variation is part of the sensitivity study; nominal resource is matched.}\label{tab:qualified_main}
\begin{tabular}{llrrrr}
\toprule ID & Role & Minimum $R$ [mm] & Minimum $\ua$ [J\,m$^{-2}$] & Maximum $\om$ & Nominal $\ua$ [J\,m$^{-2}$]\\\midrule
V00 & Sinusoidal reference & 0.96675 & 10.35180 & 0.58173 & 12.69798\\
V01 & NURBS energy extreme & 0.93598 & 11.17054 & 0.60141 & 13.68619\\
V07 & Largest-radius NURBS & 1.57747 & 9.30197 & 0.57172 & 11.42222\\
V08 & Qualified low-stress alternative & 0.99724 & 9.08890 & 0.56443 & 11.15943\\\bottomrule
\end{tabular}
\end{table*}

The varied medium resource ranges from 0.25405 to 0.28350 mm across the scenarios. This is recorded explicitly because re-adjusting thickness to restore nominal material after a manufacturing perturbation would define a different uncertainty problem. Minimum potential and maximum stress generally occur in different scenarios. Their pairing in Eq.~\eqref{eq:robust} deliberately represents two finite-set extrema, not one physically realized state.

\begin{figure*}[t]
\centering\includegraphics[width=.98\textwidth]{qualified_profiles.pdf}
\caption{Qualified examples under the same nominal initial resource and envelope. Dashed lines show reference geometry; solid lines show the mesh-24 deformed centreline colored by longitudinal extreme-fibre stress. Reaction plots show nominal response and the full 27-scenario range. V07 is the normalized finite-scenario knee within the qualified NURBS shortlist.}\label{fig:profiles}
\end{figure*}

\subsection{Numerical resolution and full evaluation cost}
The nominal search, scenario qualification, targeted refinement and fine-mesh retries together contain 764 expensive path attempts, in addition to three disclosed pilot calls. All 449 nominal and 270 scenario solver paths satisfy the inherited algorithmic checks. The initial targeted refinement and its first extension supply 21 mesh-32 paths, six mesh-40 paths and six nine-point load paths. Maximum 24--32 changes are 1.20\% in potential, 0.334\% in stress utilization and 6.35\% in reaction. The reaction sensitivity of V01 motivated the additional fine checks.

All six default mesh-64 paths reached the inherited iteration limits without satisfying the full success predicate. They remain recorded as rejected. Six retries increased the primary iteration cap from 1800 to 6000, leaving the mechanics, load protocol, tolerances and success predicate unchanged; all six then passed. No rejected high-resolution response enters the comparisons.

The sine changes by at most 0.171\% in potential between meshes 40 and 64. For V01, potential changes by 1.83--1.85\%, stress utilization by 0.65--0.89\%, and reaction by 0.58--0.77\% in absolute value over this interval. At the three checked mesh-64 scenarios, its potential advantage over the matched sine remains positive, 13.11--14.33\%, with 4.87--6.07\% greater stress utilization. These are checks at the original nominal and extremal scenarios, not a repeated 27-scenario scan at mesh 64. The sign of the interpreted trade-off persists, while its magnitude remains resolution-dependent. The 7.91\% and 3.38\% results therefore specifically describe the mesh-24 finite-scenario calculation, rather than converged continuum values.

Increasing the nominal load path from five to nine points changes terminal potential by at most 0.000098\%, stress utilization by 0.0253\% and reaction by 0.290\% in the six checks. Continuous resource matching is exact by construction; the largest polygonal-mesh resource deficit is 0.758\% at mesh 16 and 0.335\% at mesh 24. Thus small differences near the nominal front must be read alongside mesh sensitivity. The finer evidence, failed attempts and state arrays are retained in the supplement and numerical package.

